% Aşama 1 ara raporu: veri hazırlığı, bölme düzenleri, modeller ve ilk bulgular.
% Derleme: pdflatex asama1_rapor.tex (iki kez). Gerekli paketler: babel-turkish, booktabs, geometry, hyperref.
\documentclass[11pt,a4paper]{article}

\usepackage[utf8]{inputenc}
\usepackage[T1]{fontenc}
\usepackage[turkish,shorthands=off]{babel}
\usepackage{lmodern}
\usepackage[margin=2.3cm]{geometry}
\usepackage{amsmath}
\usepackage{booktabs}
\usepackage{array}
\usepackage{caption}
\usepackage{float}
\usepackage{enumitem}
\usepackage[hidelinks]{hyperref}

\setlist{itemsep=2pt, topsep=4pt}
\captionsetup{font=small, labelfont=bf}
\newcommand{\Rsq}{R$^2$}
\newcommand{\best}[1]{\textbf{#1}}

\title{Sektörler Arası Risk Yayılımı\\[2pt]\large Aşama 1 Ara Raporu: Veri Hazırlığı, Kıyas Modelleri ve Klasik LSTM}
\author{} % yazar adı eklenebilir
\date{22 Kasım 2026}

\begin{document}
\maketitle

\begin{abstract}
\noindent
Fama--French 49 sektör portföyünün günlük getirileri kullanılarak, her sektörün gelecek dönemdeki riski
(log gerçekleşen varyans) tahmin edildi. Sektörün yalnızca kendi geçmişini kullanan modeller, diğer
sektörlerin bilgisini de kullanan modellerle karşılaştırıldı. Aylık görünümde iki bölme düzeninde de
(tek bölme ve walk-forward) aynı sonuç çıktı: sektör başına ağırlık paylaşan LSTM, klasik HAR modeliyle
istatistiksel olarak berabere kaldı. Bütün sektörleri tek bir girdi olarak gören LSTM ise aşırı öğrendi ve
her yılda HAR'ın gerisinde kaldı. Piyasa genelinin risk bilgisini ekleyen HAR+piyasa modeli günlük veride
HAR'ı anlamlı biçimde geçti. Bu modelin en büyük kazancı stres dönemlerinde görüldü. Günlük LSTM koşuları
henüz yapılmadı (GPU kotası bekleniyor).
\end{abstract}

% =====================================================================
\section{Amaç}
Araştırma sorusu: bir sektörde başlayan piyasa şoku hangi koşullarda başka sektörlere yayılıyor, ve
sonraki işlem günlerinde hangi sektörlerde risk artıyor? Aşama~1'in amacı ortak bir veri hattı ve
değerlendirme düzeni kurmak, güçlü kıyas modellerini belirlemek ve klasik bir LSTM ile şu soruyu
cevaplamaktır: \emph{diğer sektörlerin geçmişi, bir sektörün gelecekteki riskini tahmin etmeye yardım ediyor mu?}

% =====================================================================
\section{Veri}
\begin{itemize}
  \item \textbf{Kaynak:} Kenneth R.\ French Data Library, \emph{49 Industry Portfolios [Daily]} (CRSP, 202609 sürümü).
        NYSE, AMEX ve NASDAQ hisselerinden oluşan değer ağırlıklı portföyler; temettü dahil yüzde getiriler, 2 ondalık.
  \item \textbf{Dönem:} 1985-01-02 $\to$ 2026-09-30; 10\,517 işlem günü, 501 ay. Bu dönemde eksik değer yok.
  \item \textbf{Sektörler:} 49 sektör veri setindeki ad ve sırasıyla kullanıldı. Hiçbir sektör birleştirilmedi,
        bölünmedi ya da gruplanmadı.
  \item \textbf{Yuvarlama:} 2\,188 hücre (\%0.4) tam olarak 0.00. Bunlar eksik veri değil, 2 ondalığa yuvarlamanın
        sonucu. Log dönüşümlerinde $\varepsilon = 0.01\ (\%^2)$ eklendi.
\end{itemize}

\paragraph{Hedefin seçimi.} Ön analizde getirinin kendisi bir gün sonrası için tahmin edilemez çıktı:
bir günlük gecikmeli korelasyon kendi sektörüyle $-0.010$, diğer sektörlerle ortalama $-0.016$.
Getirinin büyüklüğü ($|r|$) ise belirgin biçimde tahmin edilebilir: kendi sektörüyle $0.258$, diğer sektörlerle $0.203$.
Ayrıca piyasa genelindeki bir şoktan sonraki gün diğer sektörlerin riski normalin 1.6--1.75 katına çıkarken,
yalnızca tek bir sektöre özgü şoklarda bu oran 0.9--1.2 kat. Bu nedenle hedef olarak \textbf{ileri risk} seçildi.

% =====================================================================
\section{Görünümler ve veri hazırlığı}
Aynı model, günlük ya da aylık görünümden gelen veriyle beslenebilecek şekilde katmanlı bir yapı kuruldu:
kaynak $\to$ görünüm $\to$ ortak panel $\to$ bölme, ölçekleme ve pencereleme $\to$ model.
Model, verinin frekansını bilmez; $(B, L, N, F)$ boyutunda girdi alır ve $(B, N)$ tahmin üretir.

\begin{table}[h]
\centering\small
\caption{İki görünüm. $r_t$: günlük yüzde getiri; RV$_k$: son $k$ günün $r^2$ ortalaması.}
\label{tab:views}
\begin{tabular}{@{}p{2.6cm}p{6.2cm}p{6.2cm}@{}}
\toprule
 & \textbf{Günlük} & \textbf{Aylık} (danışman önerisi) \\
\midrule
Dönem & işlem günü (10\,517) & takvim ayı (501) \\
Özellikler & $r_t$,\ $\log(r_t^2+\varepsilon)$,\ $\log(\text{RV}_5+\varepsilon)$,\ $\log(\text{RV}_{22}+\varepsilon)$ &
             ayın ortalama, en yüksek ve en düşük günlük getirisi \\
Hedef $y_t$ & $\log\big(\overline{r^2}_{t+1\ldots t+5}+\varepsilon\big)$: sonraki 5 günün riski &
              $\log\big(\overline{r^2}_{\text{sonraki ay}}+\varepsilon\big)$: sonraki ayın riski \\
Girdi penceresi & 60 gün $\Rightarrow$ $(60, 49, 4)$ & 12 ay $\Rightarrow$ $(12, 49, 3)$ \\
\bottomrule
\end{tabular}
\end{table}

\noindent
Girdiler her (sektör, özellik) çifti için yalnızca train verisinin ortalama ve standart sapmasıyla standartlaştırıldı
ve $\pm10$'da kırpıldı. Hedef de sektör bazında standartlaştırıldı; tahminler değerlendirmeden önce log-RV
ölçeğine geri çevrildi. Özelliklerin yalnızca $t$ ve öncesine baktığı birim testleriyle doğrulandı.

% =====================================================================
\section{Train / val / test bölmesi}
Danışmanın 6 yıllık testi uzun bulması üzerine üç düzen tanımlandı (Tablo~\ref{tab:splits}). Hepsinde iki koruma var:
\begin{itemize}
  \item \textbf{Purge:} Bir örneğin hedef penceresi başka bir bölüme (örneğin train'den val'e) taşamaz.
  \item \textbf{Embargo:} Train örnekleri val/test dönemlerine en az 22 işlem günü (günlük) ya da 1 ay (aylık) uzak kalır.
\end{itemize}
Model ve ölçekleyici her katmanda yalnızca o katmanın train verisiyle eğitilir. Val, erken durdurma ve Ridge'in ceza
katsayısını seçmek için kullanılır.

\begin{table}[h]
\centering\small
\caption{Bölme düzenleri ve kullanılabilir örnek sayıları (train / val / test).}
\label{tab:splits}
\begin{tabular}{@{}llp{6.6cm}r@{}}
\toprule
\textbf{Düzen} & \textbf{Görünüm} & \textbf{Dönemler} & \textbf{Örnek sayısı} \\
\midrule
B (tek bölme) & günlük & train $\le$ 2010, val 2011--2016, test 2017-01 $\to$ 2026-09 & 6\,451 / 1\,505 / 2\,444 \\
B (tek bölme) & aylık  & aynı dönemler & 299 / 71 / 116 ay \\
A (walk-forward) & aylık & 10 katman. Her test yılı $Y$ (2017--2026) için train $\le Y{-}7$, val $Y{-}6 \ldots Y{-}1$,
                         test $Y$. Train her yıl 12 ay büyür. Toplam test, B ile aynı 116 ay. & 299--407 / 71 / 12 \\
Kriz pencereleri & günlük & COVID 2020-02-20 $\to$ 06-30 (train $\le$ 2018, val 2019) & 8\,463 / 252 / 92 \\
                 &        & SVB 2023-03-01 $\to$ 05-31 (train $\le$ 2021, val 2022) & 9\,220 / 251 / 64 \\
                 &        & Son 1 yıl 2025-10 $\to$ 2026-09 (train $\le$ 2024-09) & 9\,909 / 245 / 246 \\
\bottomrule
\end{tabular}
\end{table}

% =====================================================================
\section{Modeller}
\begin{table}[H]
\centering\small
\caption{Modeller. Bütün kıyas modelleri iki görünümde de aynı kodla çalışır; günlük HAR pencereleri 1, 5 ve 22 gün,
aylık HAR pencereleri 1, 3 ve 12 aydır.}
\label{tab:models}
\begin{tabular}{@{}lp{12.4cm}@{}}
\toprule
Persistence & Gelecek dönemin riski = son dönemin riski (günlükte son 5 gün, aylıkta bu ay). \\
EWMA & Üstel ağırlıklı varyans, $\lambda = 0.94$ (günlük) / $0.97$ (aylık). \\
HAR-kendi & Corsi (2009). Her sektör için ayrı en küçük kareler regresyonu: sektörün üç farklı penceredeki log riski. \\
HAR+piyasa & HAR-kendi'ye aynı üç değişkenin 49 sektör üzerinden ortalaması eklenir (piyasa göstergesi; gruplama değildir). \\
Ridge-tüm & Her sektör için 49 sektörün bütün HAR değişkenleri (147 girdi). Ceza katsayısı val'de 13 değerlik ızgaradan seçilir. \\
\midrule
LSTM-çapraz & Bütün sektörlerin geçmişini birlikte görür: $(L, 49{\cdot}F)$ $\to$ LSTM $\to$ dropout $\to$ doğrusal katman (49 çıktı).
              Gizli boyut 64 (günlük) / 16 (aylık), dropout 0.2 / 0.3. \\
LSTM-kendi & Her sektör yalnızca kendi geçmişini görür; aynı LSTM ağırlıkları paylaşılarak her sektöre ayrı uygulanır.
             Sektörü ayırt etmek için 8 / 4 boyutlu sektör gömmesi eklenir. Yapı gereği sektörler arası etki sıfırdır. \\
\bottomrule
\end{tabular}
\end{table}

\noindent
\textbf{Eğitim:} Standartlaştırılmış hedefte MSE kaybı; AdamW (öğrenme oranı $10^{-3}$, ağırlık azaltma $10^{-4}$ günlük /
$10^{-3}$ aylık), gradyan kırpma 1.0, val kaybına göre erken durdurma (sabır 10 / 20, en fazla 100 / 300 epoch).
Her model 5 farklı seed ile eğitildi; raporlanan tahmin, 5 seed'in ortalamasıdır.

\noindent
\textbf{Neden iki LSTM?} İki model arasındaki tek fark, diğer sektörlerin geçmişini görüp görmemeleridir.
Bu nedenle aralarındaki fark, Granger nedenselliğinin doğrusal olmayan bir karşılığı olarak okunabilir.

% =====================================================================
\section{Değerlendirme ölçütleri}
\begin{itemize}
  \item \textbf{OOS \Rsq{}} (ana ölçüt): $1 - \text{SSE}_{\text{model}} / \text{SSE}_{\text{train ortalaması}}$, sektör bazında
        hesaplanıp 49 sektörün ortalaması alınır. 0 değeri ``hep geçmiş ortalamayı söylemek kadar iyi'' demektir.
        Walk-forward'da her tahmin, o tarihte bilinen train ortalamasıyla kıyaslanır.
  \item \textbf{Göreli MSE:} Modelin MSE'si / HAR-kendi'nin MSE'si. 1'in altı HAR'dan iyi demektir.
  \item \textbf{QLIKE:} Varyans ölçeğinde Patton (2011) kaybı. Riskin eksik tahminini fazla tahmininden daha ağır cezalandırır.
  \item \textbf{Diebold--Mariano (DM) testi:} Her dönemin 49 sektör ortalamalı kaybı üzerinden yapılır. Newey--West gecikmesi
        $h-1$ (günlük 4, aylık 0) ve Harvey--Leybourne--Newbold düzeltmesi kullanıldı. İstatistik negatifse ilk model daha iyidir.
\end{itemize}
Bir modeli ``başarılı'' saymak için ölçüt: HAR-kendi'yi göreli MSE $<1$ ve DM $p<0.05$ ile geçmesi.

% =====================================================================
\section{Bulgular}

\subsection{Aylık görünüm: LSTM sonuçları}

\begin{table}[h]
\centering\small
\caption{Aylık test sonuçları (2017-01 $\to$ 2026-08, 116 ay). DM sütunu: HAR-kendi'ye karşı istatistik ($p$).
A için değerler 10 walk-forward katmanının birleşik tahminlerinden hesaplandı.}
\label{tab:monthly}
\begin{tabular}{@{}lrrrrrrr@{}}
\toprule
 & \multicolumn{3}{c}{\textbf{B (tek bölme)}} & & \multicolumn{3}{c}{\textbf{A (walk-forward)}} \\
\cmidrule{2-4}\cmidrule{6-8}
\textbf{Model} & OOS \Rsq & Göreli MSE & DM ($p$) & & OOS \Rsq & Göreli MSE & DM ($p$) \\
\midrule
Persistence  & 0.273 & 1.236 & 2.59 (0.010) & & 0.267 & 1.241 & 2.65 (0.008) \\
EWMA         & 0.094 & 1.568 & 5.95 ($<$0.001) & & 0.085 & 1.575 & 5.96 ($<$0.001) \\
HAR-kendi    & 0.415 & 1.000 & -- & & 0.413 & 1.000 & -- \\
HAR+piyasa   & \best{0.424} & \best{0.990} & $-0.70$ (0.482) & & \best{0.424} & \best{0.986} & $-1.00$ (0.317) \\
Ridge-tüm    & 0.349 & 1.154 & 6.54 ($<$0.001) & & 0.349 & 1.151 & 6.84 ($<$0.001) \\
LSTM-çapraz  & 0.291 & 1.298 & 5.42 ($<$0.001) & & 0.301 & 1.277 & 5.07 ($<$0.001) \\
LSTM-kendi   & 0.417 & 1.005 & 0.18 (0.859) & & 0.411 & 1.010 & 0.34 (0.733) \\
\midrule
\multicolumn{8}{@{}l}{LSTM-çapraz -- LSTM-kendi: DM $= 7.92$ ($p<0.001$) [B]; \ $6.83$ ($p<0.001$) [A]} \\
\bottomrule
\end{tabular}
\end{table}

\begin{table}[h]
\centering\small
\caption{A düzeni, yıl bazında HAR-kendi'ye göre göreli MSE (her yıl 12 ay; 2026: 8 ay). 1'in altı HAR'dan iyi.}
\label{tab:yearly}
\begin{tabular}{@{}lrrrrrrrrrr@{}}
\toprule
\textbf{Model} & 2017 & 2018 & 2019 & 2020 & 2021 & 2022 & 2023 & 2024 & 2025 & 2026 \\
\midrule
HAR+piyasa  & 0.99 & 0.98 & 1.01 & 0.97 & 1.04 & 0.94 & 0.96 & 1.04 & 1.00 & 0.99 \\
LSTM-kendi  & 1.22 & 1.06 & 0.89 & 0.99 & 1.10 & 0.98 & 0.93 & 0.97 & 0.98 & 1.17 \\
LSTM-çapraz & 1.73 & 1.21 & 1.01 & 1.19 & 1.43 & 1.56 & 1.32 & 1.26 & 1.08 & 2.00 \\
\bottomrule
\end{tabular}
\end{table}

\begin{enumerate}
  \item \textbf{LSTM-kendi, HAR-kendi ile berabere.} Fark iki düzende de anlamlı değil ($p = 0.86$ ve $0.73$).
        LSTM-kendi 10 yılın 6'sında HAR'ı geçti, 4'ünde geride kaldı. Seed'ler arası fark çok küçük
        (val kaybı 0.456--0.460).
  \item \textbf{LSTM-çapraz aşırı öğreniyor.} Train kaybı 0.35'e inerken val kaybı yaklaşık 40. epoch'ta dibe vurup yükseliyor.
        Seed'ler arası fark büyük (val kaybı 0.537--0.617). Model 10 yılın hepsinde HAR'dan, iki düzende de LSTM-kendi'den
        anlamlı derecede kötü. Doğrusal tarafta da aynı tablo var: Ridge-tüm, HAR-kendi'nin gerisinde.
        Yani 147 girdiyi 299--407 aylık örnekle hiçbir yapı vermeden öğrenmek mümkün olmuyor.
  \item \textbf{A ile B aynı sonucu veriyor.} Modelleri her yıl yeni veriyle yeniden eğitmek sıralamayı değiştirmedi.
        Aylık bulgu tek bir bölmeye bağlı değil. Ayrıca y2017 katmanının eğitim kayıtları B koşusuyla birebir aynı;
        bu, koşuların tekrarlanabilir olduğunu gösteriyor.
\end{enumerate}

\subsection{Kıyas modelleri: günlük görünüm}

\begin{table}[h]
\centering\small
\caption{Günlük kıyas modelleri: OOS \Rsq{} ve HAR-kendi'ye göre göreli MSE. ``Birleşik'': üç kriz/son-yıl penceresi birlikte.}
\label{tab:daily}
\begin{tabular}{@{}lrrcrrrr@{}}
\toprule
 & \multicolumn{2}{c}{\textbf{B (2017--2026)}} & & \multicolumn{4}{c}{\textbf{Göreli MSE, kriz pencereleri}} \\
\cmidrule{2-3}\cmidrule{5-8}
\textbf{Model} & OOS \Rsq & Göreli MSE & & COVID & SVB & Son 1 yıl & Birleşik \\
\midrule
Persistence & 0.013 & 1.523 & & 0.878 & 1.601 & 1.678 & 1.374 \\
EWMA        & 0.283 & 1.107 & & 1.128 & 1.030 & 1.050 & 1.076 \\
HAR-kendi   & 0.354 & 1.000 & & 1.000 & 1.000 & 1.000 & 1.000 \\
HAR+piyasa  & \best{0.367} & \best{0.984} & & \best{0.857} & 0.985 & \best{0.988} & \best{0.940} \\
Ridge-tüm   & 0.338 & 1.038 & & 0.904 & \best{0.973} & 1.117 & 1.019 \\
\bottomrule
\end{tabular}\\[4pt]
\begin{minipage}{0.86\linewidth}\footnotesize
HAR+piyasa -- HAR-kendi DM: B'de $-2.21$ ($p=0.027$; sektörlerin tek tek testinde 49 sektörün 26'sında anlamlı iyileşme);
kriz pencereleri birleşik $-1.75$ ($p=0.080$). Kriz pencerelerindeki LSTM koşuları GPU kotası nedeniyle henüz yapılmadı.
\end{minipage}
\end{table}

\begin{enumerate}[resume]
  \item \textbf{Piyasa bilgisi günlük veride yardım ediyor.} HAR+piyasa, HAR-kendi'yi B'de anlamlı biçimde geçti.
        Kazanç COVID penceresinde en büyük: MSE \%14 daha düşük. Aylık veride de HAR+piyasa üç stres döneminin üçünde
        HAR-kendi'den iyi (COVID 0.94, 2022 0.94, SVB 0.91), sakin dönemlerde ise avantajı yok. Bu, piyasa geneline yayılan
        şoklarda diğer sektörlerin durumunun tahmin için önem kazandığını gösteren bir ipucu; ancak aylık kesitlerde
        örnek sayısı çok az (5, 12 ve 3 ay).
  \item \textbf{Kriz pencerelerinde OOS \Rsq{} yanıltıcı.} COVID penceresinde bütün modellerin \Rsq{} değeri 0.68--0.77.
        Bunun nedeni modellerin başarısı değil: risk train ortalamasının çok üstüne çıktığı için, riskin yükseldiğini fark eden
        her model yüksek \Rsq{} alıyor. Kriz pencerelerinde göreli MSE ve DM esas alınmalı.
  \item \textbf{QLIKE'ta EWMA iyi görünüyor.} QLIKE riskin eksik tahminini daha ağır cezalandırıyor, EWMA ise yapısı gereği riski
        biraz yüksek tahmin ediyor. Bu yüzden ana ölçüt olarak log ölçeğinde MSE ve \Rsq{} kullanıldı.
\end{enumerate}

% =====================================================================
\section{Sınırlar}
\begin{itemize}
  \item Aylık test 116 aydan oluşuyor; yıl bazında ve kriz kesitlerinde örnek sayısı çok az. 10 ayrı yıllık testte
        1--2 tanesinin şans eseri anlamlı çıkması beklenir; ana sonuç birleşik değerlendirmedir.
  \item LSTM'ler için hiperparametre araması yapılmadı; her görünümde tek bir ayar denendi.
  \item Günlük LSTM sonuçları henüz yok; günlükte yalnızca kıyas modelleri raporlandı.
  \item Yalnızca değer ağırlıklı portföyler ve tek bir risk vekili (log ortalama kare getiri) kullanıldı.
  \item Aylık toplama ay içindeki zamanlamayı yok ediyor: şok ve tepki çoğu zaman aynı aya düşüyor.
\end{itemize}

% =====================================================================
\section{Sonraki adımlar}
\begin{enumerate}
  \item Günlük kriz pencerelerinde LSTM koşusu (\texttt{configs/lstm\_daily\_crisis.yaml}, GPU gerekiyor).
  \item LSTM-çapraz için düzenlileştirme denemesi: daha küçük model, daha yüksek dropout ve ağırlık azaltma ile val üzerinde
        küçük bir ızgara araması. Amaç, kaybın yalnızca kötü ayardan kaynaklanıp kaynaklanmadığını görmek.
  \item Şok senaryosu analizi: bir sektörün girdisine $-k\sigma$ şok verip diğer sektörlerin tahminindeki değişimden
        49$\times$49 yayılım matrisi çıkarmak; sakin ve gergin piyasa rejimlerinde karşılaştırmak.
  \item Aşama 2--4: GARCH ve VAR tabanlı yayılım endeksleri, ardından sektörler arası etkileşimi bir graf üzerinden modelleyen
        GNN / Temporal GNN. Aylık bulgu da buna işaret ediyor: ağırlık paylaşan model çalışıyor, yapısız yoğun model çöküyor.
\end{enumerate}

% =====================================================================
\section*{Tekrarlanabilirlik}
\small\sloppy
Kod: \url{https://github.com/salihcengiz/sector-spillover} (commit \texttt{89c0e91} ve sonrası). Deney ayarları
\texttt{configs/} altında: \texttt{lstm\_monthly\_b.yaml}, \texttt{lstm\_monthly\_a.yaml}, \texttt{lstm\_daily\_crisis.yaml},
\texttt{lstm\_daily\_b.yaml}. LSTM koşuları Google Colab'da (CPU) yapıldı; çıktılar Drive'da \texttt{20261122-064031\_lstm\_monthly\_b}
ve \texttt{20261122-071641\_lstm\_monthly\_a} klasörlerinde. Günlük kıyas modeli sonuçları aynı kodla (LSTM'siz, birkaç saniyelik
hesap) üretildi. Bütün akış 47 birim testiyle sınanıyor: özelliklerin geleceği görmemesi, purge ve embargo, ölçekleyicinin yalnızca
train'e bakması ve sektör modelinde sektörler arası etkinin sıfır olması.

\end{document}
